\chapter{机器学习建模、调参与模型评估}

\section{实验设置与评价指标}

本章对应课程评分标准中的“机器学习建模、调参与模型评估”。全部实验使用 Python、pandas、numpy、scipy、matplotlib 和 scikit-learn，在 CPU 环境下完成，不使用深度学习。随机过程统一设置 \texttt{random\_state=42}。监督学习任务使用训练/测试划分，训练阶段进行交叉验证和参数搜索，测试集只用于最终评估。

分类任务评价指标包括 Accuracy、Precision、Recall、F1、混淆矩阵，并补充 ROC-AUC、PR-AUC 和 Balanced Accuracy。高风险筛查场景中，Recall 和 F1 比单纯 Accuracy 更重要。回归任务评价指标包括 $R^2$、MSE、MAE，并补充 RMSE。聚类任务使用肘部法则和 Silhouette，并补充 Calinski-Harabasz 和 Davies-Bouldin。

\section{分类任务：高风险数字生活方式筛查}

分类目标变量为 \texttt{high\_risk\_flag}。模型包括 Logistic Regression、Random Forest 和 Gradient Boosting。训练阶段使用 StratifiedKFold 交叉验证进行调参，随后在验证集上进行阈值选择。由于高风险筛查不能只追求 Accuracy，本文重点比较默认阈值、最大 F1 阈值、Recall 至少 0.60 和 Recall 至少 0.70 的策略。

表 \ref{tab:classification-threshold-tuning} 给出验证集阈值调优结果，表 \ref{tab:classification-tuned-metrics} 给出测试集指标。最终采用 Gradient Boosting + threshold=0.14，在测试集上 Recall=0.6420、F1=0.5355、PR-AUC=0.5084。该策略属于偏召回的筛查策略，不能解释为医学诊断工具。

\begin{table}[H]
\centering
\caption{验证集阈值调优结果}
\label{tab:classification-threshold-tuning}
\resizebox{\textwidth}{!}{
\begin{tabular}{llrrrrrr}
\toprule
模型 & 阈值策略 & 阈值 & Precision & Recall & F1 & PR-AUC & Balanced Accuracy \\
\midrule
gradient\_boosting & default\_0\_50 & 0.5000 & 0.6429 & 0.2045 & 0.3103 & 0.4805 & 0.5880 \\
gradient\_boosting & max\_f1 & 0.3500 & 0.5752 & 0.4924 & 0.5306 & 0.4805 & 0.7005 \\
gradient\_boosting & recall\_at\_least\_60\_best\_precision & 0.1400 & 0.4175 & 0.6136 & 0.4969 & 0.4805 & 0.6992 \\
gradient\_boosting & recall\_at\_least\_70\_best\_precision & 0.1200 & 0.2636 & 0.7727 & 0.3931 & 0.4805 & 0.6149 \\
\bottomrule
\end{tabular}
}
\end{table}


\input{tables/classification_tuned_metrics}

图 \ref{fig:classification-metrics} 比较不同阈值策略的分类指标。默认阈值 0.50 的 Precision 较高但 Recall 较低；threshold=0.14 能识别更多高风险样本，但误报也会增加，这符合筛查模型在召回和精确率之间的权衡。

\maybefigure{classification_tuned_metrics_comparison.png}{分类阈值策略指标对比}{fig:classification-metrics}

图 \ref{fig:classification-confusion} 为最终阈值下的混淆矩阵。测试集中 $TN=566$、$FP=133$、$FN=63$、$TP=113$，说明模型能够召回一部分高风险样本，但仍然需要人工复核和谨慎解释。

\maybefigure{classification_final_confusion_matrix.png}{最终分类混淆矩阵}{fig:classification-confusion}

图 \ref{fig:classification-pr} 和图 \ref{fig:classification-roc} 分别展示 PR 曲线和 ROC 曲线。在高风险正类较少的情况下，PR-AUC 对正类识别能力更敏感；ROC-AUC 则用于补充观察整体排序能力。

\maybefigure{classification_precision_recall_curve.png}{分类 Precision-Recall 曲线}{fig:classification-pr}

\maybefigure{classification_roc_curve.png}{分类 ROC 曲线}{fig:classification-roc}

\section{回归任务：数字依赖预测与生产力弱预测}

回归任务以 \texttt{digital\_dependence\_score} 为主线，同时保留 \texttt{productivity\_score} 作为对照。表 \ref{tab:regression-target-comparison} 显示，数字依赖预测效果明显强于生产力预测。Gradient Boosting 在数字依赖目标上的测试集 $R^2=0.9839$、MSE=3.1471、MAE=0.9982；而生产力得分预测 $R^2=-0.0041$，说明当前特征对该目标解释能力弱。

\begin{table}[H]
\centering
\caption{两个回归目标的最佳模型对比}
\label{tab:regression-target-comparison}
\resizebox{\linewidth}{!}{
\begin{tabular}{llrrrrr}
\toprule
目标变量 & 最佳模型 & MAE & MSE & RMSE & $R^2$ & CV 最优 $R^2$ \\
\midrule
\texttt{productivity\_score} & Gradient Boosting & 7.1671 & 85.2031 & 9.2306 & -0.0041 & 0.0119 \\
\texttt{digital\_dependence\_score} & Gradient Boosting & 0.9982 & 3.1471 & 1.7740 & 0.9839 & 0.9804 \\
\bottomrule
\end{tabular}
}
\end{table}


\begin{table}[H]
\centering
\caption{\texttt{digital\_dependence\_score} 回归模型测试集指标}
\label{tab:regression-digital-dependence}
\begin{tabular}{lrrrr}
\toprule
模型 & MAE & MSE & RMSE & $R^2$ \\
\midrule
Gradient Boosting & 0.9982 & 3.1471 & 1.7740 & 0.9839 \\
Ridge & 0.7652 & 3.6541 & 1.9116 & 0.9813 \\
Linear Regression & 0.7651 & 3.6538 & 1.9115 & 0.9813 \\
Random Forest & 1.2680 & 4.3875 & 2.0946 & 0.9776 \\
\bottomrule
\end{tabular}
\end{table}


图 \ref{fig:regression-target-comparison} 展示两个回归目标的指标差异，图 \ref{fig:regression-dd-observed} 展示数字依赖真实值与预测值对比。数字依赖预测中真实值与预测值接近对角线，说明数字行为变量对数字依赖分数具有较强表征能力。

\maybefigure{regression_target_comparison.png}{两个回归目标预测效果对比}{fig:regression-target-comparison}

\maybefigure{regression_digital_dependence_observed_vs_predicted.png}{数字依赖得分真实值与预测值对比}{fig:regression-dd-observed}

图 \ref{fig:regression-dd-importance} 展示数字依赖预测的置换重要性。重要变量主要包括 \texttt{device\_hours\_per\_day}、\texttt{notifications\_per\_day} 和 \texttt{phone\_unlocks}，与数字依赖指标的含义一致。但强预测关系不等于因果关系，合成数据结果也不能直接外推到现实场景。

\maybefigure{regression_digital_dependence_permutation_importance.png}{数字依赖预测置换重要性}{fig:regression-dd-importance}

生产力弱预测结果见表 \ref{tab:regression-productivity}。负结果说明，同一套数字行为和生活习惯特征并不能稳定解释所有 outcome，报告中不能写成“数字生活方式可以有效预测生产力”。

\begin{table}[H]
\centering
\caption{\texttt{productivity\_score} 回归模型测试集指标}
\label{tab:regression-productivity-metrics}
\begin{tabular}{lrrrrr}
\toprule
模型 & CV 最优 $R^2$ & MAE & MSE & RMSE & $R^2$ \\
\midrule
gradient\_boosting & 0.0119 & 7.1671 & 85.2031 & 9.2306 & -0.0041 \\
random\_forest & 0.0072 & 7.2088 & 85.9208 & 9.2693 & -0.0125 \\
ridge & 0.0043 & 7.1747 & 85.5884 & 9.2514 & -0.0086 \\
linear\_regression & 0.0002 & 7.1999 & 85.8232 & 9.2641 & -0.0114 \\
\bottomrule
\end{tabular}
\end{table}


\section{聚类任务：数字生活方式用户画像}

聚类任务只使用数字行为与生活习惯数值特征，不使用人口背景类别和 outcome 字段训练。模型比较包括 KMeans、AgglomerativeClustering 和 GaussianMixture。表 \ref{tab:clustering-model-comparison} 显示，KMeans 在 $k=3$ 时 Silhouette=0.1860，结合可解释性作为最终画像方案。

\begin{table}[H]
\centering
\caption{聚类算法与簇数对比摘要}
\label{tab:clustering-model-comparison}
\begin{tabular}{lrrrr}
\toprule
算法 & $k$ & Silhouette & Calinski-Harabasz & Davies-Bouldin \\
\midrule
KMeans & 2 & 0.1808 & 752.8023 & 1.9065 \\
KMeans & 3 & 0.1860 & 611.1770 & 1.7959 \\
KMeans & 4 & 0.1406 & 546.4364 & 1.9615 \\
Agglomerative & 3 & 0.1518 & 495.8938 & 1.8768 \\
GaussianMixture & 2 & 0.1478 & 650.8461 & 2.1599 \\
GaussianMixture & 4 & 0.1225 & 408.1725 & 2.2266 \\
\bottomrule
\end{tabular}
\end{table}


图 \ref{fig:clustering-elbow} 和图 \ref{fig:clustering-silhouette} 分别展示 KMeans 肘部曲线和不同算法/簇数的 Silhouette 对比。Silhouette 整体较低，说明分群边界偏弱，因此聚类仅用于探索性画像，不代表天然清晰人群边界。

\maybefigure{clustering_kmeans_elbow.png}{KMeans 肘部曲线}{fig:clustering-elbow}

\maybefigure{clustering_silhouette_by_k.png}{不同聚类算法与簇数的 Silhouette 对比}{fig:clustering-silhouette}

图 \ref{fig:pca-explained-variance} 展示 PCA 累计解释方差，图 \ref{fig:clustering-pca} 展示聚类结果在前两个 PCA 维度上的投影。前两个主成分累计解释约 42.41\% 方差，PCA 仅用于可视化和辅助理解。

\maybefigure{pca_explained_variance.png}{PCA 累计解释方差曲线}{fig:pca-explained-variance}

\maybefigure{clustering_lifestyle_pca.png}{数字生活方式聚类的 PCA 可视化}{fig:clustering-pca}

表 \ref{tab:clustering-profiles-compact} 和图 \ref{fig:clustering-heatmap} 总结三类画像：高社交媒体使用型、高设备依赖型和低负荷均衡型。这些名称只用于描述簇均值特征，不是固定人群标签。

\begin{table}[H]
\centering
\caption{数字生活方式聚类画像精简表}
\label{tab:clustering-profiles-compact}
\resizebox{\textwidth}{!}{
\begin{tabular}{rrrrrrrrrl}
\toprule
簇 & 样本数 & 设备时长 & 社交媒体分钟 & 睡眠时长 & 睡眠质量 & 高风险比例 & 生产力均值 & 数字依赖均值 & 建议画像名 \\
\midrule
0 & 447 & 6.9533 & 430.2125 & 7.3816 & 2.7899 & 0.1969 & 65.3375 & 34.3296 & 高社交媒体使用型 \\
1 & 1039 & 11.0538 & 131.8787 & 6.1218 & 1.6953 & 0.3725 & 65.9081 & 51.2425 & 高设备依赖型 \\
2 & 2014 & 5.4711 & 113.4275 & 7.8106 & 3.2137 & 0.1142 & 64.9767 & 29.6962 & 低负荷均衡型 \\
\bottomrule
\end{tabular}
}
\end{table}


\maybefigure{clustering_lifestyle_profile_heatmap.png}{聚类画像标准化特征热力图}{fig:clustering-heatmap}

\section{模型结果综合分析}

综合三类任务，分类模型可作为高风险初筛，但需要人工复核；数字依赖回归效果较好，是本文回归主线；生产力预测为弱预测/负结果，说明当前特征不足以稳定解释该目标；聚类画像具有解释价值，但 Silhouette 较低，边界不强。

\section{业务场景解释与结果边界}

本文结果可用于数字生活方式管理参考，例如关注设备使用时长、通知干扰、社交媒体使用、睡眠和身体活动平衡。分类结果只适合初筛，不应用于医学诊断、自动化惩罚或强制干预。由于数据为合成数据，所有结论都应作为课程建模和行为分析结果谨慎解释。
